% Define document class
\documentclass[twocolumn]{aastex631}
\usepackage{showyourwork}

\usepackage{amsmath}
\usepackage{amssymb}
\usepackage{xspace}
\usepackage{mathtools}

\DeclarePairedDelimiter\abs{\lvert}{\rvert}
\DeclarePairedDelimiter\norm{\lVert}{\rVert}

% Swap the definition of \abs* and \norm*, so that \abs
% and \norm resizes the size of the brackets, and the 
% starred version does not.
\makeatletter
\let\oldabs\abs
\def\abs{\@ifstar{\oldabs}{\oldabs*}}
%
\let\oldnorm\norm
\def\norm{\@ifstar{\oldnorm}{\oldnorm*}}
\makeatother

\newcommand{\comment}[1]{}

% note command for inline notes
\newcommand{\note}[1]{\textsf{\textcolor{red}{#1}}}
% command to note a missing citation
\newcommand{\needscite}{\note{[needs citation]}}

\newcommand{\lya}{Lyman-$\alpha$\xspace}
\newcommand{\du}{\Delta u}
\newcommand{\overbar}[1]{\mkern 5mu\overline{\mkern-5mu#1\mkern-5mu}\mkern 5mu}



\shorttitle{Modeling Galaxy Catalogues with Normalizing Flows}
\shortauthors{Crenshaw et al.}


% Begin!
\begin{document}

\title{Feasibility of Detecting a Photometric Lyman-alpha Signal with the Vera Rubin Observatory}

\correspondingauthor{John Franklin Crenshaw}
\email{jfc20@uw.edu}

\author[0000-0002-2495-3514]{John Franklin Crenshaw}
\affiliation{DIRAC Institute and Department of Physics, University of Washington, Seattle, WA 98195, USA}

\author{Anze Slosar} % Orcid?
\affiliation{}

\author{Alex Malz} % Orcid?
\affiliation{}

\author{Marilena Loverde} % Orcid?
\affiliation{}

\author{Sam Schmidt} % Orcid?
\affiliation{}

\author{Melissa Graham} % Orcid?
\affiliation{}

\author{Eric Gawiser} % Orcid?
\affiliation{}

\author{Jeff Newman} % Orcid?
\affiliation{}

\author{Daniel Gruen} % Orcid?
\affiliation{}

\author{Andrew Connolly} % Orcid?
\affiliation{}

\collaboration{100}{(Dark Energy Science Collaboration)}

% Abstract with filler text
\begin{abstract}
    The scattering of light by the \lya transition of neutral hydrogen in the intergalactic medium produces the so-called ``\lya Forest'', visible in the observed spectra of distant quasars.
    This same effect should decrease the flux of distant sources observed by the Vera Rubin Observatory.
    Considering only statistical errors and cosmic variance, we assess the feasibility of detecting this signal with Rubin photometry via a decrease in $u$ band flux for sources at redshift $> 2.36$.
    We estimate a signal-to-noise ratio (SNR) of NN for the mean \lya absorption, NN for the \lya autocorrelation, NN for the galaxy clustering cross-correlation, NN for the lensing cross-correlation, and NN for the dust cross-correlation.
    We estimate that including infrared photometry from the Euclid and Roman space telescopes will increase SNR by a factor of NN.
    Given the high SNR estimated here, we expect this signal to be systematics limited, and discuss the modeling that must be performed to enable a robust detection with Rubin data.
    We also discuss how this photometric \lya probe is complementary to spectroscopic \lya probes and the standard Rubin 3x2pt cosmology analysis.
\end{abstract}

% Main body with filler text
\section{Introduction}
\label{sec:intro}

\note{
    The space between galaxies is filled with a diffuse gas that accounts for most of the baryonic content of the universe \needscite.
    Most of this gas is hydrogen, and after reionization, most of this hydrogen is ionized.
    However, a small fraction of this hydrogen is neutral (describe why - how is it protected).
    Despite it's low density, this gas has a very high cross section with photons at $\lambda_\alpha$, which excites the \lya transition.
    This causes the neutral gas to scatter light at this wavelength out of the line-of-sight of the observer.
    Many such scatterings along the line-of-sight create a plethora of absorption lines blueward of the \lya emission line.
    These many absorption lines can be observed in the spectra of high-redshift objects.
    This feature has been named the \lya Forest.
}

\note{
    The neutral hydrogen gas traces the large scale structure of the diffuse baryons in the IGM.
    This allows us to use the \lya Forest as a comsological tracer of large scale structure.
    Spectroscopic surveys use statistics of the \lya Forest to measure properties of the gas, and measure the clustering of matter in the universe.
    Maybe also Dark Energy?
    Constrain Dark Matter?
    Many citations here.
}

\note{
    Photometric surveys like the Rubin observatory do not have the spectral resolution to identify the trees, but they may be able to see the forest.
    High-redshift galaxies will have their blue light scattered by the \lya Forest, which will result in observing lower magnitudes in blue bands than would otherwise be expected.
    Furthermore, this effect will be greater for galaxies that are behind overdensities, and less for galaxies that are behind underdensities.
}

\note{
    Detecting this effect relies on selecting high-redshift galaxies, predicting how bright they should be in the blue bands, and comparing to the observed magnitudes.
    This is hard, as predicting redshifts from photometric data is noisy and prone to systematic errors.
    In addition, bluer bands typically have lower SNR to begin with.
    And the effect is relatively small.
    However, the hope is that when averaged over many sources in a huge survey like LSST, the signal will rise above the noise and we will be able to see the forest for the trees.
}

\note{
    What photometric surveys lack in spectral resolution, they make up for in statistics.
    This enables us to do complementary things.
    We are insensitive to longitudinal modes, but can better constrain the transverse modes.
    In addition, we can directly measure the mean extinction of the \lya Forest, which is something that must be marginalized over in spectroscopic studies.
    Providing a direct measurement will allow spectroscopic studies to tighten their constraints.
}

\note{
    In this paper, we make optimistic assumptions (i.e. no systematic errors), and attempt to forecast how well the Rubin observatory could measure the \lya and calculate correlation functions.
    We consider auto- and cross-correlations, and calculate the signal to noise.
    Probably say something about the redshift samples as well.
    We consider how adding NIR data could improve the analysis.
    Finally, we point to the potential of this method, and discuss the sytematic errors that need to be modeled in order to bring this method to maturity.
    We also discuss this as a potential systematic for high-redshift photo-z's that may impact other cosmology analyses.
    Hopefully, we can test the method during early LSST science.
}

\section{The photometric Lyman-alpha signal}

\label{sec:signal}

Assume that a distant galaxy has an intrinsic SED $f_0(\lambda)$ in the rest frame of the observer.
\lya scattering in the IGM modifies the SED, resulting in the observed SED $f(\lambda) = F(\lambda) f_0(\lambda)$, where $F(\lambda)$ is the transmission of the \lya Forest.
Note that $F(\lambda)$ traces the density of neutral hydrogen gas at redshift $z = \lambda / \lambda_\alpha - 1$ along the line-of-sight.

Photometric surveys do not observe the SED directly, but rather the flux of the SED integrated over a bandpass:
\begin{align}
    f = \int d\lambda \, R(\lambda) F(\lambda) f_0(\lambda),
\end{align}
where $R(\lambda)$ is the response function of the photometric band.
We can define the \emph{photometric transmission}:
\begin{align}
    F = \frac{f}{f_0} = \int d\lambda \, R(\lambda) F(\lambda) s(\lambda),
    \label{eq:F}
\end{align}
where $s(\lambda)$ is the shape of the SED, defined
\begin{align}
    s(\lambda) = \frac{f_0(\lambda)}{f_0} = \frac{f_0(\lambda)}{\int d\lambda \, R(\lambda) f_0(\lambda)}.
\end{align}

The rest frame wavelength of the \lya transition is $\lambda_\alpha = 1216$\AA, and the Rubin $u$ band has support over the wavelengths $3200 < \lambda < 4085$.
This means that galaxies at $z > 1.63$ will appear dimmer than expected in the $u$ band because some of this light will be scattered by neutral hydrogen along the line-of-sight. 
This effect will be greatest for galaxies at $z > 2.36$ whose entire $u$ band will be impacted by this scattering.
Higher redshift galaxies will appear dimmer in longer wavelength bands as well.

With a few assumptions, we can estimate the impact  of \lya scattering on LSST $u$ band observations of galaxies at redshifts $z > 2.36$.
We will assume a fiducial SED shape $\bar{s}(\lambda) = 1$ over the support of the $u$ band.
The expected transmission can be calculated from the effective optical depth of the \lya Forest reported by [https://arxiv.org/pdf/1904.01110.pdf]: $\bar{F}(\lambda) = e^{-\tau_\text{eff}(\lambda)}$, where
\begin{align}
    \tau_\text{eff} = \tau_0 (1 + z)^\gamma,
\end{align}
$\tau_\text{eff} = (5.54 \pm 0.64) \times 10^{-3}$ and $\gamma = 3.182 \pm 0.074$.
With these choices, we can calculate the expected size of the photometric signal in LSST:
\begin{align}
    \bar{F} = \int d\lambda \, R(\lambda) \bar{F}(\lambda) \bar{s}(\lambda) = ...
    \label{eq:Fbar}
\end{align}
This translates to an expected $u$ band decrement $\overbar{\du} = -2.5 \log_{10} \bar{F} = ...$.
In other words, galaxies at $z > 2.36$ will, on average, appear $...$ magnitudes dimmer in the $u$ band than would otherwise be expected in the absence of the \lya Forest.

The question is whether this decrease in flux can be detected with the Rubin observatory and what cosmological signal can be extracted from it.
In Section \ref{sec:snr}, we simulate observations and the inference pipeline to calibrate the expected SNR of the signal on a per-galaxy basis.
In Section \ref{sec:correlations}, we calculate auto- and cross-correlations of the photometric \lya signal to determine the expected SNR of the signal in a cosmological analysis.


\section{Estimating signal-to-noise}
\label{sec:snr}

In this section, we simulate Rubin observations to calibrate the expected SNR of the photometric \lya signal on a per-galaxy basis.
We will not simulate variations in the \lya signal due to large scale structure, because these variations are greatly subdominant to the uncertainties in estimating $\du$, and will therefore not impact our signal-to-noise estimate in this section.
For an estimate of the SNR on measurements of the \lya variations due to large scale structure, see Section \ref{sec:correlations}

\subsection{The simulated catalog}

We use the simulated galaxy catalog of [Cite both of Melissa's papers] \needscite, which is based on the Millennium simulation \needscite.
This catalog contains \note{N} galaxies up to redshift \note{N}, with optical LSST photometry and near-infrared (NIR) photometry from Euclid and the Roman Space Telescope.
This constitutes the truth catalog.

Throughout the rest of this paper, we will denote true galaxy redshifts\footnote{Occasionally, $z$ will denote the LSST $z$ band, but this will be obvious from context.} with the variable $z$, true LSST $u$ band magnitudes with $u$, and LSST $u$ band magnitudes including \lya extinction with the variable $u_\alpha$.
We will use the variable $m$ to refer to the vector of all other magnitudes.
For most of the paper, $m$ refers only to the LSST $grizy$ bands.
In Section \ref{sec:nir}, we consider the inclusion of Euclid $YJH$ and Roman $YJHF$ bands.

To simulate an observed catalog, we add $u$ band decrements to galaxies at redshifts $z > 1.63$ to account for the \lya scattering of the $u$ band flux.
We calculate these decrements by integrating the mean \lya transmission up to the wavelength of the $u$ band filter that is impacted by the \lya transmission:
\begin{align}
    \du = -2.5 \log_{10} \left[ \, 
    \int_0^{\lambda_z} d\lambda \, R(\lambda) \bar{F}(\lambda)
    \, \right],
\end{align}
where $\lambda_z = \min(\lambda_\alpha (1 + z), 2.36)$.

We then simulate photometric errors according to the \needscite error model, including the low SNR generalization presented in \needscite and the extension to Euclid and Roman presented in \needscite.
We denote the observed magnitudes with variables $\hat{u}_\alpha$, $\hat{m}$, and their photometric errors with the variables $\sigma_u$, $\sigma_m$.
Finally, we apply a cut of SNR > 5 in the $i$ band.

\note{
    At some point I want to add plots showing how different assumptions of $\bar{s}$ change the mean $\du$.
    Make a plot $\du(z)$ for several different powerlaws for $\bar{s}$.
}

\subsection{Selecting the high-redshift sample}
\label{sec:photo-z}

We train an ensemble of normalizing flows to model the conditional distribution $p(z, u | m)$ \needscite.
Note this ensemble is trained on the truth catalog, imagining a galaxy simulation that matches reality.
\note{
    I need to actually write this section.
    I can probably refer the ensemble setup to that described in my pzflow paper.
}

Once we have trained the model, we create a high-redshift sample whose $u$ band magnitudes are maximally impacted by \lya scattering via applying the photometric redshift cut
\begin{align}
    \int_{2.36}^\infty dz ~ p(z|\hat{m}, \sigma_m) > 0.95.
    \label{eq:redshift-cut}
\end{align}
That is, we select those galaxies whose redshift posteriors indicate a 95\% probability of being above redshift 2.36.
We can efficiently compute the integral in Equation \ref{eq:redshift-cut} by sampling
\begin{align}
    _{i=1}^{~~ M} \{ m^i \, | \, m^i \sim p(m | \hat{m}, \sigma_m) \},
\end{align}
from the photometric error model for each galaxy (which is just a Gaussian in flux space), followed by sampling
\begin{align}
    _{j=1}^{~~ N} \{ z^{(i,\, j)}, u^{(i,\, j)} \, | \, z^{(i,\, j)}, u^{(i,\, j)} \sim p(z, u | m^i) \},
    \label{eq:samples}
\end{align}
from the flow ensemble, and calculating the fraction of $\{ z^{(i, j)} \}$ samples that are above 2.36.

\note{
    State the values I use for $M, N$.
    Maybe also leave $p_\text{cut}$ general, and state my choice when I state $M, N$?
}

\note{
    Discuss the completeness and purity of the redshift sample.
    Discuss its growth as a function of survey time.
    Discuss how the probability cut was selected to favor purity, and show the plot of the tradeoff between purity and completeness.
}

\note{
    Project the size of the sample for LSST.
}

\subsection{Inference on $\du$}

Now that we have a sample of galaxies at redshift > 2.36, we can perform inference on the $u$ band decrement $\Delta u$.
This amounts to using the observed $grizy$ magnitudes to predict $u$, and comparing this predicted value to the observed value, $\hat{u}_\alpha$.

The likelihood for the decrement of each galaxy can be calculated:
\begin{align}
    \mathcal{L}(\du) = \int du \, p(\hat{u}_\alpha | u, \du) \, p(u | \hat{m}, z > 2.36).
\end{align}
This integral can be efficiently computed via reusing the samples from Equation \ref{eq:samples}.
For the galaxies in the high-redshift sample, we reject samples with $z^{(i, j)} < 2.36$, generating a set of samples 
\begin{align}
    {}_{i=1}^{~~ M} {}_{j=1}^{~~ N} \{ u^{(i,\, j)} | z^{(i,\, j)} > 2.36\}
\end{align}
for each galaxy.
We can then directly calculate $p(\hat{u}_\alpha | u, \du) = p(\hat{u}_\alpha | u_\alpha = u + \du)$ using the photometric error model.
Averaging over the samples $\{ u^{(i,\, j)} \}$ yields the likelihood $\mathcal{L}(\du)$.

You can see likelihoods for three example galaxies in Figure \note{N}.
\note{
    Note how they are non-Gaussian, and try to include one that has a large tail due to low signal to noise.
    Note that the likelihoods span to negative decrements - so you can't even tell if the galaxy actually did get dimmer - it may be brighter!
}

\note{
    When combining many likelihoods, the variance scales like $\bar{\sigma}_{\small \du} / N$, where $\bar{\sigma}_{\small \du}$ is the effective Gaussian error per galaxy.
    Note that the effective SNR for a single galaxy is $\sim 1$.
    This isn't great SNR, but in the next section we will calculate correlation functions that have much greater SNRs.
}


\section{Auto- and cross-correlations of the Lyman-alpha signal}

The \lya Forest traces neutral hydrogen gas in the universe, which in turn traces the large scale structure of the universe.
We can calculate correlations of the \lya signal with itself and other tracers of the large scale structure to learn about the clustering and behavior of the different components of the universe.
This also allows us to combine information from many galaxies and boost the relatively modest SNR of the photometric \lya signal for a single galaxy, as well as test for systematic errors that may plague our analysis.
I will first develop the formalism of the photometric \lya tracer, before calculating auto- and cross-correlations in the following subsections.

\lya correlation functions are usually computed in terms of the transmission $F(\lambda)$, which traces the density of neutral hydrogen gas at redshift $z = \lambda / \lambda_\alpha - 1$ along the line-of-sight.
For the rest of this section, we will work in terms of redshift $z$ instead of wavelength $\lambda$. 
This requires us to make the replacement $d\lambda \to \lambda_\alpha dz$ in our integrals, where $\lambda_\alpha$ is the wavelength of the \lya transition.

We can write the transmission in terms of fluctuations about the mean: $F(z) = \bar{F}(z) [1 + \delta_F(z)]$.
Similarly, we can express the shape of each galaxy's continuum SED over the $u$ band in terms of fluctuations about the mean SED: $s(z) = \bar{s(z)}[1 + \delta_s(z)]$.
Note that the mean $\bar{F}$ is an average over lines-of-sight, while the mean $\bar{s}(z)$ is an average over galaxies in the sample.

Using Equations \ref{eq:F} and \ref{eq:Fbar}, we can calculate fractional variations in the \lya transmission $F$:
\begin{align}
    \Delta_F = \frac{F - \bar{F}}{\bar{F}} = \int dz \, W_F(z)[\delta_F(z) + \delta_s(z)],
\end{align}
where $W_F(z) \propto R(z) \bar{F}(z) \bar{s}(z)$ is a normalized redshift kernel for the photometric \lya signal.
$\delta_F$ represents the \lya signal, while $\delta_s$ represents noise from galaxy-to-galaxy variations of the continuum SED.
The noise $\delta_s$ should not correlate with the signal \note{[?]}, and we will assume that the continuum SED can be estimated well enough that this term will be subdominant.
Thus, we have the tracer
\begin{align}
    \Delta_F = \int dz \, W_F(z) \delta_F(z).
    \label{eq:lya-tracer}
\end{align}

Equation \ref{eq:lya-tracer} is the photometric \lya tracer.
We see that it is a projection of variations in the \lya transmission field, $\delta_F(z)$, weighted by the kernel $W_F(z)$.
In Fourier space, 
\begin{align}
    \delta_F(\textbf{k}) = b_F (1 + \mu_\textbf{k}^2 \beta_F) \delta_b(\textbf{k}),
    \label{eq:fourier-tracer}
\end{align}
where $b_F \approx -0.13$ is the linear bias of the \lya field, $\beta_F \approx 1.5$ measures the strength of redshift-space distortions (RSD)\footnote{Unlike the galaxy tracer (c.f. Equation \ref{eq:galaxy-tracer}), $\beta$ is another free parameter. \note{This has to do with the nonlinear transformation applied to the optical depth? I need to work on understanding this.}}, $\mu_\textbf{k}$ is the cosine between the mode $\textbf{k}$ and the line-of-sight, and $\delta_b(\textbf{k})$ are the Fourier space fluctuations of the baryon density field \needscite.
Note that the bias is negative as more neutral hydrogen corresponds to less transmitted flux.

The bias and RSD make this tracer look like the projected galaxy density tracer (c.f. Equation \ref{eq:galaxy-tracer}), but similar to gravitational lensing tracers (c.f. Section \note{N}), it is sensitive to density integrated over the lines-of-sight to background sources.
The integral in Equation \ref{eq:lya-tracer} means the photometric tracer is relatively insensitive to the RSD term in Equation \ref{eq:fourier-tracer}.
This is because the RSD term only contributes to modes along the line-of-sight (due to the cosine $\mu_\textbf{k}$).
Modes with wavelength less than the width of the redshift kernel will cancel each other in the integral, and modes longer than the width of the redshift kernel will be indistinguishable from \note{$\bar{F}$ or something to do with $s$?}.
\note{
    I should test if the correlations are relatively independent of $\beta$!
    I should really have Anze confirm this part.    
}

In the subsections that follow, I calculate different correlations of the photometric \lya signal and forecast the SNR of these measurements as a function of LSST survey duration.
In every case, correlation functions are calculated using CCL \needscite assuming the Planck cosmology \needscite.
The SNR calculations involve the expected per-galaxy uncertainty on the photometric \lya transmission, $\bar{\sigma}_{\Delta_F}$, which can be calculated from the expected per-galaxy uncertainty on the $u$ band decrement, $\bar{\sigma}_{\small \du}$:
\begin{align}
    \bar{\sigma}_{\Delta_F} = \frac{\ln 10}{2.5} \bar{\sigma}_{\small \du} \approx ... 
\end{align}
\note{Say something about the logarithmic binning of tracers that applies to every subsection.}

\subsection{Lyman-alpha autocorrelation}

\note{
    Write this after I actually perform the calculation. 
    Explain the counting of pairs.
    Don't forget to add the factor of 1/2 to account for double counting.
}

\subsection{Correlation with galaxy clustering}

We can correlate the photometric \lya signal with foreground galaxies that lie within the same redshift range as the HI clouds responsible for the \lya scattering of the $u$ band flux.
In order to perform this cross-correlation, we must photometrically select a sample of foreground tracer galaxies.
This selection can be performed in an identical manner to the photo-z selection in Section \ref{sec:photo-z} by replacing the limits on the integral in Equation \ref{eq:redshift-cut}:
\begin{align}
    \int_{1.63}^{2.36} dz ~ p(z|\hat{m}, \sigma_m) > 0.95.
\end{align}
Again, this cut yields a high-purity but low-completeness sample.
The purity, completeness, and projected size of the foreground sample can be seen as a function of LSST survey duration in Figure \note{N}.

The tracer for the foreground galaxies can be written:
\begin{align}
    \delta_g(\textbf{k}) = b_g (1 + \mu_\textbf{k}^2 \beta_g) \delta_b(\textbf{k}),
    \label{eq:galaxy-tracer}
\end{align}
where $b_g \approx 2$ is the linear galaxy bias and $\beta_g \approx b_g^{-1}$ sets the amplitude of RSDs \needscite.
There is a corresponding projected galaxy tracer
\begin{align}
    \Delta_g = \int dz \, W_g(z) \, \delta_g(z),
\end{align}
where $W_g(z)$ is the redshift distribution of galaxies in the foreground sample\footnote{This distribution is typically written as $n(z)$, but I write it as $W_g(z)$ here to unify the notation with the redshift kernels of the other tracers considered in this section.}

\note{Plot the correlations.}

\note{Give the pair counting that yields the SNR estimate.}


\subsection{Correlation with lensing magnification}

\note{
    Need to ask Marilena/Anze about whether I should use the lensing magnification signal in the number density or in the magnitude increase.
    If doing the latter, I should use $i$ band, as this will be the highest SNR band in my catalog.
}

\subsection{Correlation with dust}

\note{
    Can calculate correlation with galaxy colors as a proxy for dust.
    Anze thinks this will be interesting to CMB people.
}

\note{
    Look at Marilena's notes in the email she sent me.
    Eric Gawiser might be a good person to ask.
    Or even Daniel Gruen?
}


\section{Impact of near-infrared photometry}



\end{document}
